\documentclass[11pt]{article}
\usepackage{amsmath, amssymb}
\usepackage{geometry}
\usepackage{algorithm}
\usepackage{algpseudocode}
\usepackage{graphicx}
\usepackage{hyperref}

\geometry{margin=1in}

\title{Mathematical Description of BATPING Pulse Optimization}
\author{BATPING Project}
\date{May 26, 2026}

\begin{document}
\maketitle

\tableofcontents
\newpage

%==============================================================================
\section{Introduction and Motivation}
%==============================================================================

\subsection{The BATPING System}

BATPING is an acoustic temperature measurement system designed for telescope dome seeing prediction. The fundamental principle is straightforward: the speed of sound in air depends on temperature according to:
\begin{equation}
    v_s = 331.3 \sqrt{1 + \frac{T_C}{273.15}} \approx 331.3 + 0.606 \, T_C \quad \text{[m/s]}
\end{equation}
where $T_C$ is the temperature in Celsius. By precisely measuring the time-of-flight (TOF) of an acoustic pulse over a known distance $d$, we can infer the air temperature:
\begin{equation}
    T_C = \frac{1}{0.606}\left(\frac{d}{\Delta t} - 331.3\right)
\end{equation}

For a path length of $d = 18.7\,\text{cm}$ and speed of sound $v_s \approx 343\,\text{m/s}$ (at 20°C), the acoustic travel time is approximately:
\begin{equation}
    \Delta t = \frac{d}{v_s} = \frac{0.187}{343} \approx 545\,\mu\text{s}
\end{equation}

\subsection{Why Pulse Shape Matters}

The precision of TOF measurement depends critically on how sharply we can identify the arrival time of the acoustic pulse. Consider two extreme cases:

\begin{enumerate}
    \item \textbf{Infinitely narrow pulse:} An ideal Dirac delta function $\delta(t - t_{\text{arrival}})$ would provide perfect timing precision, but carries zero energy and is physically unrealizable.

    \item \textbf{Continuous wave:} A sustained sinusoid carries maximum energy but provides no timing information---every cycle looks identical.
\end{enumerate}

Real pulses must balance these extremes. We seek a pulse that is:
\begin{itemize}
    \item \textbf{High amplitude:} Strong signal relative to noise allows reliable detection
    \item \textbf{Narrow width:} Sharp features enable precise arrival time identification
\end{itemize}

\subsection{The Piezoelectric Transducer Challenge}

Our system uses a 40\,kHz piezoelectric transducer array. Piezoelectric transducers are \textit{resonant} devices---they vibrate most efficiently at their natural frequency (40\,kHz in our case). This resonance is characterized by the quality factor $Q$:
\begin{equation}
    Q = \frac{f_0}{\Delta f}
\end{equation}
where $f_0$ is the resonant frequency and $\Delta f$ is the bandwidth.

High-$Q$ resonators present a fundamental problem: they exhibit \textbf{ringdown}. When excited, they continue oscillating long after the driving signal stops, with amplitude decaying as:
\begin{equation}
    A(t) \propto e^{-\pi f_0 t / Q}
\end{equation}

This ringdown ``smears out'' our pulse in time, degrading TOF precision. The optimization problem we address is: \textit{Can we shape the electrical input waveform to minimize the acoustic output's temporal spread while maintaining adequate amplitude?}

\subsection{Our Approach: Black-Box Inverse Filtering}

In principle, if we knew the transducer's transfer function $H(f)$, we could compute a pre-distorted input:
\begin{equation}
    X_{\text{input}}(f) = \frac{X_{\text{desired}}(f)}{H(f)}
\end{equation}
such that the output $Y(f) = H(f) \cdot X_{\text{input}}(f) = X_{\text{desired}}(f)$ matches our desired pulse shape. This is classical \textbf{inverse filtering}.

However, direct inverse filtering is impractical because:
\begin{enumerate}
    \item The transfer function $H(f)$ is difficult to measure precisely in situ
    \item Noise amplification occurs where $|H(f)|$ is small
    \item The system may have nonlinearities not captured by a linear transfer function
\end{enumerate}

Instead, we take a \textbf{black-box optimization} approach: we parameterize the input waveform, measure the actual acoustic output, and iteratively adjust the parameters to improve the output. This is equivalent to searching for the inverse filter without explicitly computing it---letting the optimization algorithm discover the pre-distortion empirically.

%==============================================================================
\section{Waveform Parameterization}
%==============================================================================

\subsection{The Carrier Signal}

Since our transducer resonates at 40\,kHz, we must drive it near this frequency for efficient energy transfer. The carrier signal is simply:
\begin{equation}
    c(t) = \sin(2\pi f_c t)
\end{equation}
where $f_c = 40\,\text{kHz}$.

\subsection{The Envelope Function}

Rather than transmitting a continuous 40\,kHz tone, we modulate its amplitude with an \textit{envelope function} $e(t)$ that shapes when and how strongly the carrier is present:
\begin{equation}
    x(t) = e(t) \cdot \sin(2\pi f_c t)
\end{equation}

The envelope determines the pulse's temporal structure. Common choices include:
\begin{itemize}
    \item \textbf{Rectangular:} $e(t) = 1$ for $t \in [0, T_p]$, zero otherwise. Simple but causes spectral splatter.
    \item \textbf{Gaussian:} $e(t) = \exp\left(-\frac{(t-t_0)^2}{2\sigma^2}\right)$. Smooth and spectrally compact.
    \item \textbf{Hanning/Hamming:} Raised cosine windows. Good spectral properties.
\end{itemize}

\subsection{Fourier Series Parameterization}

To enable optimization, we need a \textit{parameterized} family of envelope functions. We use a truncated Fourier series, which can approximate any periodic function:
\begin{equation}
    e(t; \boldsymbol{\theta}) = a_0 + \sum_{k=1}^{K} \left[ a_k \cos\left(\frac{2\pi k t}{T}\right) + b_k \sin\left(\frac{2\pi k t}{T}\right) \right]
\end{equation}

Here:
\begin{itemize}
    \item $\boldsymbol{\theta} = [a_0, a_1, b_1, a_2, b_2, \ldots, a_K, b_K]^\top \in \mathbb{R}^{2K+1}$ is the parameter vector
    \item $T$ is the waveform period (4\,ms in our implementation)
    \item $K$ is the number of harmonics (we used $K = 4$, giving 8 parameters)
    \item $a_0$ is the DC (constant) component
    \item $a_k$ are cosine coefficients (symmetric contributions)
    \item $b_k$ are sine coefficients (antisymmetric contributions)
\end{itemize}

\textbf{Why Fourier?} The Fourier basis is:
\begin{enumerate}
    \item \textbf{Complete:} Can represent any square-integrable function as $K \to \infty$
    \item \textbf{Orthogonal:} Parameters are independent; changing $a_3$ doesn't affect the $a_1$ component
    \item \textbf{Interpretable:} Each coefficient controls a specific frequency component
    \item \textbf{Smooth:} Truncated Fourier series are inherently smooth (no discontinuities)
\end{enumerate}

\subsection{Non-Negativity Constraint}

Physical envelopes must be non-negative (you cannot have negative amplitude). We enforce this with a simple clipping operation:
\begin{equation}
    \tilde{e}(t; \boldsymbol{\theta}) = \max\left(e(t; \boldsymbol{\theta}), 0\right)
\end{equation}

This is a ``soft'' constraint---the optimizer can explore parameter spaces where $e(t)$ goes negative, but the actual waveform used will have those regions clipped to zero.

\subsection{Normalization}

The arbitrary waveform generator (AWG) accepts waveforms normalized to $[-1, 1]$. The final uploaded waveform is:
\begin{equation}
    \hat{x}(t; \boldsymbol{\theta}) = \frac{\tilde{e}(t; \boldsymbol{\theta}) \cdot \sin(2\pi f_c t)}{\max_t \left| \tilde{e}(t; \boldsymbol{\theta}) \cdot \sin(2\pi f_c t) \right|}
\end{equation}

The AWG then scales this by the commanded amplitude (5\,V peak-to-peak in our experiments).

\subsection{Discretization}

The AWG operates on discrete samples. We generate $N = 4000$ points at an effective sample rate of $f_s = 1\,\text{MSa/s}$:
\begin{equation}
    x_n = \hat{x}(t_n; \boldsymbol{\theta}), \quad t_n = \frac{n}{f_s}, \quad n = 0, 1, \ldots, N-1
\end{equation}

The waveform duration is therefore $T = N/f_s = 4\,\text{ms}$, and at a repetition rate of 250\,Hz (period = 4\,ms), the waveform repeats continuously with no gap between repetitions.

%==============================================================================
\section{Signal Acquisition and Processing}
%==============================================================================

\subsection{The Measurement Chain}

When we command a waveform $x(t)$, the following physical processes occur:

\begin{enumerate}
    \item \textbf{D/A Conversion:} The AWG converts digital samples to an analog voltage signal
    \item \textbf{Electrical-to-Mechanical:} The piezoelectric transducer converts voltage to mechanical vibration
    \item \textbf{Acoustic Propagation:} Sound waves travel through air at speed $v_s \approx 343\,\text{m/s}$
    \item \textbf{Mechanical-to-Electrical:} The microphone converts acoustic pressure to voltage
    \item \textbf{A/D Conversion:} The recording system digitizes the microphone signal
\end{enumerate}

We can model the entire chain as a linear time-invariant (LTI) system with impulse response $h(t)$:
\begin{equation}
    y(t) = (h * x)(t) + n(t) = \int_{-\infty}^{\infty} h(\tau) x(t - \tau) \, d\tau + n(t)
\end{equation}
where $y(t)$ is the recorded signal and $n(t)$ is additive noise.

\subsection{The Butterworth Bandpass Filter}

The raw microphone signal contains our 40\,kHz signal plus noise at all frequencies. To isolate our signal, we apply a \textbf{bandpass filter} that passes frequencies in a narrow band around 40\,kHz while attenuating everything else.

\subsubsection{Why Butterworth?}

The Butterworth filter is characterized by a \textbf{maximally flat} passband---it has no ripples in the frequency range it passes. Its magnitude response is:
\begin{equation}
    |H(f)|^2 = \frac{1}{1 + \left(\frac{f}{f_c}\right)^{2n}}
\end{equation}
where $n$ is the filter order and $f_c$ is the cutoff frequency.

For a bandpass filter with passband $[f_L, f_H]$, we combine a high-pass (cutoff $f_L$) and low-pass (cutoff $f_H$) response:
\begin{equation}
    |H_{\text{bp}}(f)|^2 = \frac{1}{1 + \left(\frac{f_L}{f}\right)^{2n}} \cdot \frac{1}{1 + \left(\frac{f}{f_H}\right)^{2n}}
\end{equation}

\subsubsection{Our Implementation}

We use a 4th-order Butterworth bandpass filter with:
\begin{itemize}
    \item Lower cutoff: $f_L = 38\,\text{kHz}$
    \item Upper cutoff: $f_H = 42\,\text{kHz}$
    \item Passband width: $\Delta f = 4\,\text{kHz}$
\end{itemize}

This $\pm 2\,\text{kHz}$ window around 40\,kHz:
\begin{itemize}
    \item Passes our signal (centered at 40\,kHz)
    \item Rejects 33\,kHz interference from motion sensors
    \item Rejects low-frequency noise (HVAC, voices, handling)
    \item Rejects high-frequency noise
\end{itemize}

\subsubsection{Zero-Phase Filtering}

We apply the filter using \texttt{filtfilt}, which filters the signal forward, then backward. This eliminates phase distortion:
\begin{equation}
    y_{\text{filtered}} = \text{filtfilt}(b, a, y) \quad \Rightarrow \quad \text{zero phase shift}
\end{equation}

Phase preservation is critical because we're measuring arrival \textit{times}---any filter-induced delay would bias our TOF measurements.

\subsection{Envelope Extraction via the Hilbert Transform}

After bandpass filtering, we have a 40\,kHz oscillating signal. To measure its amplitude variations over time, we extract the \textbf{envelope}---the curve tracing the peaks of the oscillation.

\subsubsection{The Analytic Signal}

The Hilbert transform $\mathcal{H}$ of a real signal $s(t)$ is defined as:
\begin{equation}
    \mathcal{H}\{s\}(t) = \frac{1}{\pi} \, \text{P.V.} \int_{-\infty}^{\infty} \frac{s(\tau)}{t - \tau} \, d\tau
\end{equation}
where P.V. denotes the Cauchy principal value.

In the frequency domain, the Hilbert transform is simply a $-90°$ phase shift:
\begin{equation}
    \mathcal{F}\{\mathcal{H}\{s\}\}(f) = -j \cdot \text{sgn}(f) \cdot S(f)
\end{equation}

The \textbf{analytic signal} combines the original signal with its Hilbert transform:
\begin{equation}
    s_a(t) = s(t) + j \cdot \mathcal{H}\{s\}(t)
\end{equation}

\subsubsection{Why This Works}

Consider a simple amplitude-modulated signal:
\begin{equation}
    s(t) = A(t) \cos(2\pi f_c t)
\end{equation}
where $A(t)$ is the slowly-varying envelope we want to extract.

Its Hilbert transform is (approximately, when $A(t)$ varies slowly compared to $f_c$):
\begin{equation}
    \mathcal{H}\{s\}(t) \approx A(t) \sin(2\pi f_c t)
\end{equation}

The analytic signal is therefore:
\begin{equation}
    s_a(t) = A(t) \cos(2\pi f_c t) + j \cdot A(t) \sin(2\pi f_c t) = A(t) e^{j 2\pi f_c t}
\end{equation}

The magnitude (envelope) is simply:
\begin{equation}
    |s_a(t)| = |A(t) e^{j 2\pi f_c t}| = |A(t)| \cdot |e^{j 2\pi f_c t}| = A(t)
\end{equation}

\subsubsection{Practical Computation}

In practice, we compute:
\begin{equation}
    A(t) = \sqrt{s(t)^2 + \mathcal{H}\{s(t)\}^2}
\end{equation}

This extracts the instantaneous amplitude regardless of the carrier's phase, giving us a smooth envelope that follows the peaks of the 40\,kHz oscillation.

%==============================================================================
\section{Metric Definitions}
%==============================================================================

\subsection{Time Gating}

Before extracting metrics, we apply a \textbf{time gate}---we only analyze the signal within a window $[t_1, t_2]$ where we expect the acoustic pulse to arrive.

For our setup:
\begin{itemize}
    \item Separation distance: $d = 18.7\,\text{cm}$
    \item Speed of sound: $v_s \approx 343\,\text{m/s}$
    \item Expected travel time: $\Delta t \approx 545\,\mu\text{s} \approx 0.55\,\text{ms}$
    \item System latency (triggering, buffers): $\sim 2\,\text{ms}$
    \item Gate window: approximately $[1, 10]\,\text{ms}$ depending on calibration
\end{itemize}

Time gating serves several purposes:
\begin{enumerate}
    \item \textbf{Reject EM coupling:} Electromagnetic interference from the AWG arrives at the speed of light ($t \approx 0$), long before the acoustic signal
    \item \textbf{Reject reflections:} Later arrivals from walls/objects are excluded
    \item \textbf{Reduce noise:} We only analyze the relevant portion of the recording
\end{enumerate}

\subsection{Peak Amplitude}

The peak amplitude is the maximum envelope value within the time gate:
\begin{equation}
    A_{\text{peak}} = \max_{t \in [t_1, t_2]} A(t)
\end{equation}

This metric captures signal strength. Higher $A_{\text{peak}}$ means:
\begin{itemize}
    \item Better signal-to-noise ratio
    \item More reliable detection
    \item Greater tolerance to interference
\end{itemize}

\subsection{Full Width at Half Maximum (FWHM)}

FWHM is a standard measure of pulse width. It's the duration over which the signal exceeds half its peak value:

\begin{enumerate}
    \item Find the peak: $t^* = \arg\max_{t \in [t_1, t_2]} A(t)$
    \item Compute half-maximum: $A_{1/2} = A_{\text{peak}} / 2$
    \item Find the rising edge: $t_{\text{rise}} = \max\{t < t^* : A(t) < A_{1/2}\}$
    \item Find the falling edge: $t_{\text{fall}} = \min\{t > t^* : A(t) < A_{1/2}\}$
    \item Compute FWHM: $\text{FWHM} = t_{\text{fall}} - t_{\text{rise}}$
\end{enumerate}

Mathematically:
\begin{equation}
    \text{FWHM} = t_{\text{fall}} - t_{\text{rise}}
\end{equation}
where:
\begin{align}
    t_{\text{rise}} &= \sup\left\{ t < t^* : A(t) \leq \frac{A_{\text{peak}}}{2} \right\} \\
    t_{\text{fall}} &= \inf\left\{ t > t^* : A(t) \leq \frac{A_{\text{peak}}}{2} \right\}
\end{align}

\textbf{Interpretation:} Smaller FWHM means a sharper pulse, which enables more precise arrival time identification.

\textbf{Limitation:} For signals with long ringdown tails (like our piezo transducer output), FWHM can be dominated by the tail rather than the sharp initial attack. Alternative metrics like rise time may be more relevant for TOF precision.

%==============================================================================
\section{The Objective Function}
%==============================================================================

\subsection{Formulation}

We want to simultaneously maximize amplitude and minimize width. These are competing objectives---a continuous wave maximizes amplitude but has infinite ``width,'' while a delta function minimizes width but has zero energy.

We combine them into a single scalar objective using a weighted sum:
\begin{equation}
    J(\boldsymbol{\theta}) = w_A \cdot A_{\text{peak}}(\boldsymbol{\theta}) - w_W \cdot \frac{\text{FWHM}(\boldsymbol{\theta})}{100\,\mu\text{s}}
\end{equation}

The term $\text{FWHM}/100\,\mu\text{s}$ normalizes the width to be dimensionless and roughly order-unity (typical FWHM values are hundreds of microseconds).

\subsection{Weight Selection}

We used weights $w_A = 1.0$ and $w_W = 0.3$ (later $w_W = 0.5$).

The ratio $w_A / w_W$ determines the tradeoff:
\begin{itemize}
    \item High $w_A / w_W$: Optimizer prioritizes amplitude, tolerates wider pulses
    \item Low $w_A / w_W$: Optimizer prioritizes narrow pulses, accepts lower amplitude
\end{itemize}

With $w_A / w_W = 1.0 / 0.3 \approx 3.3$, a 1\% increase in amplitude is worth approximately $3.3\,\mu\text{s}$ increase in FWHM.

\subsection{Maximization vs. Minimization}

Convention in optimization is to \textit{minimize} objective functions. Since we want to \textit{maximize} $J(\boldsymbol{\theta})$, we minimize its negative:
\begin{equation}
    \boldsymbol{\theta}^* = \arg\min_{\boldsymbol{\theta}} \left[ -J(\boldsymbol{\theta}) \right] = \arg\min_{\boldsymbol{\theta}} \left[ -w_A \cdot A_{\text{peak}} + w_W \cdot \frac{\text{FWHM}}{100} \right]
\end{equation}

%==============================================================================
\section{The Nelder-Mead Optimization Algorithm}
%==============================================================================

\subsection{Why Derivative-Free Optimization?}

Classical optimization methods like gradient descent require the gradient $\nabla J(\boldsymbol{\theta})$. In our problem, gradients are unavailable because:

\begin{enumerate}
    \item \textbf{Hardware in the loop:} Each evaluation requires physical measurement, not just computation. We cannot analytically differentiate ``send waveform to transducer, propagate through air, record with microphone.''

    \item \textbf{Measurement noise:} Even if we could compute gradients, noise in the measurements would make them unreliable.

    \item \textbf{Non-differentiable operations:} Our signal processing includes $\max$ operations (time gating, peak finding) which are not everywhere differentiable.
\end{enumerate}

\textbf{Derivative-free} methods optimize using only function values $J(\boldsymbol{\theta})$, not derivatives.

\subsection{The Simplex Method}

The Nelder-Mead algorithm maintains a \textbf{simplex}---a geometric shape with $n+1$ vertices in $\mathbb{R}^n$. In 2D, a simplex is a triangle; in 3D, a tetrahedron; in our 8-dimensional parameter space, an 8-simplex (9 vertices).

Each vertex $\boldsymbol{\theta}_i$ is a candidate solution with associated function value $J_i = J(\boldsymbol{\theta}_i)$.

\subsection{Algorithm Operations}

The algorithm iteratively transforms the simplex through four operations:

\subsubsection{Ordering}

Sort vertices by function value (for minimization of $-J$):
\begin{equation}
    J(\boldsymbol{\theta}_1) \leq J(\boldsymbol{\theta}_2) \leq \cdots \leq J(\boldsymbol{\theta}_{n+1})
\end{equation}
So $\boldsymbol{\theta}_1$ is the best point and $\boldsymbol{\theta}_{n+1}$ is the worst.

\subsubsection{Centroid}

Compute the centroid of all vertices except the worst:
\begin{equation}
    \bar{\boldsymbol{\theta}} = \frac{1}{n} \sum_{i=1}^{n} \boldsymbol{\theta}_i
\end{equation}

\subsubsection{Reflection}

Reflect the worst point through the centroid:
\begin{equation}
    \boldsymbol{\theta}_r = \bar{\boldsymbol{\theta}} + \alpha(\bar{\boldsymbol{\theta}} - \boldsymbol{\theta}_{n+1})
\end{equation}
where $\alpha > 0$ is the reflection coefficient (typically $\alpha = 1$).

\textbf{Intuition:} Move away from the worst point, toward the region where good points cluster.

\subsubsection{Expansion}

If $\boldsymbol{\theta}_r$ is the best point so far, try going further:
\begin{equation}
    \boldsymbol{\theta}_e = \bar{\boldsymbol{\theta}} + \gamma(\boldsymbol{\theta}_r - \bar{\boldsymbol{\theta}})
\end{equation}
where $\gamma > 1$ is the expansion coefficient (typically $\gamma = 2$).

\textbf{Intuition:} If reflecting helped a lot, maybe going further will help even more.

\subsubsection{Contraction}

If reflection didn't help much, contract toward the centroid:
\begin{equation}
    \boldsymbol{\theta}_c = \bar{\boldsymbol{\theta}} + \rho(\boldsymbol{\theta}_{n+1} - \bar{\boldsymbol{\theta}})
\end{equation}
where $0 < \rho < 1$ is the contraction coefficient (typically $\rho = 0.5$).

\textbf{Intuition:} The minimum might be inside the current simplex; shrink toward it.

\subsubsection{Shrinkage}

If contraction also fails, shrink the entire simplex toward the best point:
\begin{equation}
    \boldsymbol{\theta}_i \leftarrow \boldsymbol{\theta}_1 + \sigma(\boldsymbol{\theta}_i - \boldsymbol{\theta}_1), \quad i = 2, \ldots, n+1
\end{equation}
where $0 < \sigma < 1$ is the shrinkage coefficient (typically $\sigma = 0.5$).

\textbf{Intuition:} We're probably near a minimum; refine the search locally.

\subsection{Convergence}

The algorithm terminates when:
\begin{enumerate}
    \item Maximum iterations reached (we used 150)
    \item Simplex becomes sufficiently small: $\max_i \|\boldsymbol{\theta}_i - \boldsymbol{\theta}_1\| < \epsilon$
    \item Function values converge: $\max_i |J_i - J_1| < \delta$
\end{enumerate}

\subsection{Initialization}

We initialize with:
\begin{equation}
    \boldsymbol{\theta}_0 = [1.0, 0.3, 0, 0, 0, 0, 0, 0]^\top
\end{equation}

This corresponds to an envelope:
\begin{equation}
    e(t) = 1.0 + 0.3\cos\left(\frac{2\pi t}{T}\right)
\end{equation}

The initial simplex is constructed by perturbing each component of $\boldsymbol{\theta}_0$ to create $n$ additional vertices.

%==============================================================================
\section{Hardware-in-the-Loop Evaluation}
%==============================================================================

Each call to the objective function $J(\boldsymbol{\theta})$ involves real physical measurements:

\begin{algorithm}
\caption{Objective Function Evaluation}
\begin{algorithmic}[1]
\Procedure{Evaluate}{$\boldsymbol{\theta}$}
    \State \textbf{// Generate waveform from parameters}
    \State $e(t) \gets a_0 + \sum_k [a_k \cos(2\pi k t/T) + b_k \sin(2\pi k t/T)]$
    \State $x(t) \gets \max(e(t), 0) \cdot \sin(2\pi f_c t)$
    \State $\hat{x}(t) \gets x(t) / \max|x(t)|$ \Comment{Normalize to $[-1,1]$}
    \State
    \State \textbf{// Hardware operations}
    \State \textsc{UploadWaveformToAWG}($\hat{x}$)
    \State \textsc{SetAWGOutput}(\texttt{ON})
    \State \textsc{Wait}($t_{\text{settle}} = 20\,\text{ms}$) \Comment{Allow transients to decay}
    \State $y \gets \textsc{RecordMicrophone}(t_{\text{record}} = 50\,\text{ms})$
    \State \textsc{SetAWGOutput}(\texttt{OFF})
    \State
    \State \textbf{// Signal processing}
    \State $y_{\text{bp}} \gets \textsc{ButterworthBandpass}(y, f_L=38\,\text{kHz}, f_H=42\,\text{kHz})$
    \State $A \gets |\textsc{HilbertAnalytic}(y_{\text{bp}})|$ \Comment{Envelope}
    \State
    \State \textbf{// Metric extraction (within time gate)}
    \State $A_{\text{peak}} \gets \max_{t \in [t_1, t_2]} A(t)$
    \State $\text{FWHM} \gets \textsc{ComputeFWHM}(A, A_{\text{peak}})$
    \State
    \State \textbf{// Objective value}
    \State \Return $w_A \cdot A_{\text{peak}} - w_W \cdot \text{FWHM}/100$
\EndProcedure
\end{algorithmic}
\end{algorithm}

\textbf{Time per evaluation:} Approximately 0.1--0.2 seconds (dominated by recording time and hardware communication).

\textbf{Total evaluations:} 129--142 in our runs, taking several minutes total.

%==============================================================================
\section{Results and Critical Analysis}
%==============================================================================

\subsection{Optimized Parameters}

The optimizer converged to:
\begin{equation}
    \boldsymbol{\theta}^* \approx [1.003, 0.306, 0, 0, 0, 0, 0, 0]^\top
\end{equation}

This yields an envelope:
\begin{equation}
    e(t; \boldsymbol{\theta}^*) = 1.003 + 0.306\cos\left(\frac{2\pi t}{T}\right) \approx 1.0 + 0.3\cos\left(\frac{2\pi t}{T}\right)
\end{equation}

\subsection{Interpretation}

This envelope oscillates between:
\begin{align}
    e_{\min} &= 1.0 - 0.3 = 0.7 \\
    e_{\max} &= 1.0 + 0.3 = 1.3
\end{align}

Since $e(t) > 0$ for all $t$, the waveform is \textbf{always on}---it's a gently amplitude-modulated continuous wave, not a pulse.

\subsection{Duty Cycle Analysis}

Define the duty cycle as the fraction of time the waveform exceeds 10\% of its maximum:
\begin{equation}
    D = \frac{1}{T} \int_0^T \mathbf{1}\left[ |x(t)| > 0.1 \cdot \max|x| \right] dt
\end{equation}

For our optimized waveform: $D \approx 91.7\%$.

For comparison, the Gaussian baseline ($\sigma = 5$ cycles): $D \approx 10.8\%$.

\subsection{Why the Optimizer ``Cheated''}

The optimizer found a loophole:

\begin{enumerate}
    \item \textbf{Waveform period = Repetition period:} With $T = 4\,\text{ms}$ and repetition rate 250\,Hz (period = 4\,ms), consecutive waveforms abut with no gap.

    \item \textbf{Continuous wave maximizes resonant response:} The piezoelectric transducer is a resonant system. Sustained 40\,kHz excitation allows mechanical oscillation to build up, maximizing steady-state amplitude.

    \item \textbf{FWHM becomes meaningless:} For a continuous signal, the envelope never drops below half-maximum within any reasonable time window, so FWHM is either undefined or equals the entire recording length.

    \item \textbf{Width penalty is ineffective:} With FWHM pathologically large or undefined, the width penalty term $w_W \cdot \text{FWHM}/100$ doesn't constrain the solution.
\end{enumerate}

\subsection{Measured Reality}

When we carefully compared both waveforms under controlled conditions (Session 4, May 26, 2026):

\begin{center}
\begin{tabular}{lcc}
\hline
\textbf{Metric} & \textbf{Gaussian Baseline} & \textbf{``Optimized''} \\
\hline
Peak Amplitude & 1.366 & 1.348 \\
FWHM & 547\,$\mu$s & N/A (continuous) \\
Duty Cycle & 10.8\% & 91.7\% \\
Usable for TOF? & Yes & \textbf{No} \\
\hline
\end{tabular}
\end{center}

The peak amplitudes are essentially identical (0.99$\times$ ratio). The apparent 4--9$\times$ improvements reported during optimization runs were artifacts of:
\begin{itemize}
    \item Insufficient settling time before recording
    \item Recording timing not aligned with pulse arrivals
    \item Run-to-run measurement variability
\end{itemize}

\subsection{Diagnosis: Objective Function Loophole}

The optimization framework itself is sound---black-box search for pre-distortion waveforms is a valid approach to inverse filtering. The failure was in the \textbf{objective function formulation}, which allowed a degenerate solution:

\begin{enumerate}
    \item \textbf{No duty cycle constraint:} The objective function did not penalize waveforms that are ``always on.'' The optimizer was free to explore continuous waves.

    \item \textbf{FWHM fails for continuous signals:} When the envelope never drops below half-maximum, FWHM becomes undefined or pathologically large. The width penalty term becomes ineffective precisely when it's needed most.

    \item \textbf{Local optimum in wrong region:} The optimizer found a local optimum (continuous wave, high steady-state amplitude) that satisfies the mathematical objective but violates the implicit physical requirement of producing discrete pulses.
\end{enumerate}

The solution is not to abandon the optimization approach, but to \textbf{reformulate the objective function} to close this loophole.

%==============================================================================
\section{Reformulated Objective Function}
%==============================================================================

\subsection{Design Requirements}

The new objective function must:
\begin{enumerate}
    \item Reward high output amplitude (for SNR)
    \item Reward sharp arrival detection (for TOF precision)
    \item \textbf{Penalize high duty cycle} (enforce pulsed behavior)
\end{enumerate}

\subsection{Rise Time vs.\ FWHM}

For TOF measurements, what matters is how precisely we can identify the \textit{arrival instant}. The \textbf{10--90\% rise time} is more directly relevant than FWHM:
\begin{equation}
    t_{\text{rise}} = t_{90\%} - t_{10\%}
\end{equation}
where $t_{x\%}$ is the first time the envelope reaches $x\%$ of its peak value.

\textbf{Advantages of rise time:}
\begin{itemize}
    \item Measures attack sharpness directly
    \item Ignores ringdown tail (which doesn't affect arrival detection)
    \item Well-defined even for signals with extended tails
\end{itemize}

\subsection{Duty Cycle Constraint}

Define the duty cycle of the \textit{input} waveform:
\begin{equation}
    D(\boldsymbol{\theta}) = \frac{1}{T} \int_0^T \mathbf{1}\left[ |x(t; \boldsymbol{\theta})| > 0.1 \cdot \max|x| \right] dt
\end{equation}

This measures the fraction of time the waveform is ``active'' (above 10\% of maximum amplitude).

We impose a maximum allowed duty cycle $D_{\max} = 0.25$ (25\%), enforced via a penalty term:
\begin{equation}
    P_D(\boldsymbol{\theta}) = \left[ \max\left( D(\boldsymbol{\theta}) - D_{\max}, 0 \right) \right]^2
\end{equation}

The squared penalty is smooth (differentiable) and grows rapidly for large violations. A waveform with 91.7\% duty cycle would incur penalty $(0.917 - 0.25)^2 = 0.445$, which with weight $w_D = 10$ contributes $-4.45$ to the objective---a catastrophic penalty.

\subsection{Complete Reformulated Objective}

\begin{equation}
    \boxed{
    J'(\boldsymbol{\theta}) = w_A \cdot A_{\text{peak}} - w_R \cdot \frac{t_{\text{rise}}}{100\,\mu\text{s}} - w_D \cdot \left[ \max\left( D(\boldsymbol{\theta}) - 0.25, 0 \right) \right]^2
    }
\end{equation}

\textbf{Recommended weights:}
\begin{itemize}
    \item $w_A = 1.0$ (amplitude weight)
    \item $w_R = 0.5$ (rise time weight)
    \item $w_D = 10.0$ (duty cycle penalty weight)
\end{itemize}

\subsection{Expected Behavior}

With this reformulation:
\begin{enumerate}
    \item Continuous waveforms ($D > 0.25$) will be heavily penalized
    \item The optimizer will search within the space of pulsed waveforms
    \item The tradeoff becomes amplitude vs.\ rise time (attack sharpness)
    \item If a pre-distortion exists that reduces ringdown while maintaining amplitude, the optimizer should find it
\end{enumerate}

\subsection{Alternative: Constrain Waveform Duration}

An alternative approach is to use a shorter waveform that doesn't fill the entire repetition period:
\begin{equation}
    T_{\text{waveform}} < T_{\text{rep}} - T_{\text{silence}}
\end{equation}

For example, a 1\,ms waveform at 250\,Hz (4\,ms period) would have 3\,ms of guaranteed silence. This architecturally prevents continuous waveforms but reduces the parameter space the optimizer can explore.

\end{document}
